  

%%%%%%%%%%%%%%Empieza la presentacion%%%%%%%%%%%%%%%%%%%%%%%%%%%%%%%%%%%%

%\begin{frame}
%\frametitle{Esquema}
 %   \tableofcontents
%\end{frame}

\section{Objetivos}

\begin{frame}
\frametitle{Objetivos del tema}

\begin{itemize}
\item Aprender a distinguir desde un punto de vista matemático el comportamiento de las señales deterministas y los diferentes tipos de señales aleatorias (estacionarias, ergódicas, proceso de Markov).

\visible<2->{
\item Definir matemáticamente las propiedades más relevantes de:
}
\begin{itemize}
\normalsize
\visible<3->{
\item Variables aleatoria: distribución/densidad de probabilidad, valor medio, valor cuadrático medio, varianza, correlación
}


\visible<4->{
\item Señales aleatorias: valor medio, valor cuadrático medio, covarianza, autocorrelación y correlación cruzada, densidad espectral de potencia
}



\end{itemize}

\visible<5->{\item Procesamiento de señales aleatorias a través de las técnicas de procesamiento de señales deterministas}

\end{itemize}

\end{frame}

